% ============================================================
% Laborbuch_Gesamt.tex
% SafetySpot – alle drei Teilnehmer, eine chronologische Tabelle
% Autoren: D0rag3on (Noel Neuhoff), IIRu5hII (Milian), 0qln
% ============================================================

\ssSetDocTitle{Laborbuch}
\ssSetKunde{Bildungsministerium NRW}
\ssSetProjekt{SafetySpot}
\ssSetAutoren{D0rag3on (Noel Neuhoff) \textbar{} IIRu5hII (Milian) \textbar{} 0qln}
\ssSetStand{29.06.2026}
\maketitle
\setcounter{section}{0}

\section{Grundlage und Projektkontext}

Dieses Dokument vereint die persönlichen Laborbücher aller drei Projektteilnehmer
des Projekts \emph{SafetySpot} in einer einzigen, chronologisch sortierten Tabelle.
Einträge werden nicht nach Autor getrennt, sondern rein nach Datum einsortiert.
Der Autor ist in jeder Zeile explizit ausgewiesen.

Die Einträge von IIRu5hII ab Januar 2026 wurden sachlich aus Projektdefinition,
Mockups und Quellcode rekonstruiert, da frühe Arbeitsschritte ohne Commits erfolgten.
Alle anderen Einträge basieren direkt auf Git-Commits und GitHub-Pull-Requests im
Repository \texttt{safety-spot/SafetySpot}.

\section{Chronologische Laborbuch-Einträge (alle Autoren)}

\begin{sslabortabelle}
2026-01-12 & IIRu5hII & Recherche &

  Eintrag 01: Projektdefinition gesichtet und Ausgangslage festgehalten.
  Gefahrenaufklärung an Schulen erfolgt häufig nur mündlich und wenig interaktiv.
  Daraus wurde abgeleitet, dass die App niedrigschwellig, spielerisch und
  direkt verständlich sein muss.
  \newline\textit{Ergebnis:} Projektdefinition als fachliche Grundlage;
  Fokus auf Schüler, Lehrer und Schulen. \\

2026-01-12 & IIRu5hII & Ziel &

  Eintrag 02: Zielbild für die mobile App formuliert: Schüler sollen
  Gefahrensituationen in einer sicheren Umgebung bewerten, daraus lernen
  und durch Punkte/Rankings motiviert bleiben.
  \newline\textit{Ergebnis:} Zielbild für Frontend, Szenarien und Gamification übernommen. \\

2026-01-12 & IIRu5hII & Annahmen/Hypothese &

  Eintrag 03: Annahme: Ein spielerischer Ablauf mit kurzen Aufgaben und
  sofortigem Feedback ist für Schüler motivierender als reine Sicherheitsbelehrungen.
  \newline\textit{Ergebnis:} Beeinflusst Spielscreen mit zwei großen Antwortbuttons und Feedback. \\

2026-01-22 & IIRu5hII & Entscheidung &

  Eintrag 04: Entwicklungsumgebung festgelegt: Android Studio.
  Alternativen (IntelliJ IDEA, Qt Creator, Cordova) als weniger passend bewertet.
  \newline\textit{Ergebnis:} Android Studio als gemeinsame Arbeitsumgebung eingeplant. \\

2026-01-23 & IIRu5hII & Entscheidung &

  Eintrag 05: Programmiersprache für Android-App: Kotlin.
  Gründe: moderne Android-Entwicklung, Null-Safety, weniger Boilerplate,
  gute Integration in Android Studio.
  \newline\textit{Ergebnis:} Kotlin als Basis für die mobile App. \\

2026-01-23 & IIRu5hII & Recherche &

  Eintrag 06: Informationsmodell ausgewertet. Relevante Datenblöcke:
  Organisation/Rollen, Lerninhalte, Fortschritt/Gamification und Admin/Lizenz.
  \newline\textit{Ergebnis:} Datenmodell als Grundlage für spätere DTOs, Mockdaten und UI-Strukturen. \\

2026-01-23 & IIRu5hII & Entscheidung &

  Eintrag 07: Grundsatzentscheidung für Server-als-Wahrheit übernommen.
  Backend speichert alles; App hält nur Token, Cache und ggf.\ ausstehende Versuche.
  \newline\textit{Ergebnis:} Frühe Frontend-Version trotzdem mit Mockdaten aufgebaut. \\

2026-01-24 & IIRu5hII & Ziel &

  Eintrag 08: Projektvision analysiert. Erstes Release: lauffähige Android-App
  mit interaktiven Szenarien, Punkten/Ranking und einfachem Admin-Gedanken.
  \newline\textit{Ergebnis:} Scope auf wichtigste Schüler-Funktionen fokussiert. \\

2026-01-24 & IIRu5hII & Idee &

  Eintrag 09: Kategorien aus Vision und Mockups abgeleitet:
  Chemieraum, Werkraum, Sportunterricht, Technikraum, Straßenverkehr.
  \newline\textit{Ergebnis:} Kategorien in \texttt{sampleScenarios} und Home-Screen sichtbar. \\

2026-01-24 & IIRu5hII & Bewertung &

  Eintrag 10: Risiken gesichtet: Android-only, Leaderboard-Sozialdruck,
  zu komplexe eigene Szenarien.
  \newline\textit{Ergebnis:} UI einfach gehalten; Leaderboard motivierend statt bloßstellend. \\

26.04.2026 & D0rag3on & Versuchsaufbau &

  Eintrag 11: Initialisierung des Server-Projekts unter dem vorläufigen Namen
  \texttt{SafetySpotServer}. Maven-basiertes Spring-Boot-Projekt inkl.
  \texttt{mvnw}, \texttt{pom.xml}, \texttt{persistence.xml}, \texttt{beans.xml}.
  Zu diesem Zeitpunkt noch kein funktionaler Code – lediglich ein lauffähiger,
  leerer Startpunkt als Grundlage für die Backend-Entwicklung. \\

26.04.2026 & 0qln & Versuchsaufbau &

  Eintrag 12: Repository \texttt{safety-spot/SafetySpot} initialisiert
  (Commit \texttt{183ac5f}).
  Grundlegende Repo-Struktur: \texttt{.gitignore}, \texttt{.gitattributes},
  \texttt{README.md}; Zeilenendenrichtlinien (\texttt{.gitattributes})
  für plattformübergreifende Zusammenarbeit normalisiert.
  \newline\textit{Ergebnis:} Gemeinsames Repository als Startpunkt für das gesamte Team. \\

09.06.2026 & D0rag3on & Versuchsaufbau &

  Eintrag 13: Migration des Backend-Projekts von Maven auf Gradle sowie
  Umbenennung von \texttt{SafetySpotServer} auf \texttt{ssbackend}.
  Das gesamte Projekt verwendet Gradle als einheitliches Build-System;
  Maven war ein inkonsistenter Fremdkörper.
  316 Zeilen (mvnw, pom.xml, IDE-Konfiguration) entfernt.
  Neues Modul mit Spring-Boot-Gradle-Plugin und \texttt{SsbackendApplication.java}
  aufgebaut.
  \textit{Anmerkung:} Der zugehörige erste Commit hätte früher erfolgen sollen. \\

09.06.2026 & 0qln & Versuchsaufbau &

  Eintrag 14: Versionsdokumentation eingepflegt – v0-Dokumente (Projektvision,
  Informationsmodell, Entscheidungen, Risiken, Klassendiagramm \texttt{docs/v0})
  und v1-Mockups (\texttt{docs/v1}) dem Repository hinzugefügt
  (Commits \texttt{cf734b4}, \texttt{cf9eabf}).
  \newline\textit{Ergebnis:} Fachliche Grundlage und UX-Mockups im Repo versioniert;
  alle Teammitglieder haben Zugriff ohne separaten Dateiumlauf. \\

09.06.2026 & 0qln & Versuchsaufbau &

  Eintrag 15: \texttt{sscommon}-Kotlin-Modul initiiert sowie
  Copilot-Instruktionen und Projekt-Spezifikation angelegt (PR \#26,
  \texttt{1a90de6}).
  Inhalt: Copilot-Instruktionsdatei \texttt{.github/copilot-instructions.md},
  \texttt{sscommon/}-Gradle-Modul (geteilte Bibliothek),
  \texttt{spec/plan.md} (Backend-Implementierungsplan mit Paketen, DTOs,
  Security-Konzept), ISSUE-001 bis ISSUE-011 (Backend-Aufgaben).
  Außerdem: Dev-Umgebung als Nix-Flake vorbereitet.
  \newline\textit{Ergebnis:} Strukturierte Aufgabenplanung und Copilot-Wissen
  als Grundlage für alle weiteren Backend-Entwicklungsschritte. \\

09.06.2026 & 0qln & Versuchsaufbau &

  Eintrag 16: \texttt{ssmobile}-Android-Projekt initialisiert (PR \#27,
  Commit \texttt{6caaf64}).
  Gradle-Build-Konfiguration, \texttt{AndroidManifest.xml},
  Basis-Theme (Material3), Basis-Navigation und \texttt{MainActivity.kt}.
  Struktur angelegt für spätere UI-Screens.
  \newline\textit{Ergebnis:} Lauffähiges Android-Projekt-Skelett;
  IIRu5hII kann sofort auf der gleichen Basis aufbauen. \\

2026-06-10 & IIRu5hII & Ziel &

  Eintrag 17: Start der praktischen Frontend-Arbeit ohne Backend-Integration.
  Ziel: klickbarer Android-Prototyp, der Mockups nachbildet und den Kernablauf
  von Login bis Szenario-Spiel zeigt.
  \newline\textit{Ergebnis:} Rekonstruiert aus Codezustand und Autorenaussage. \\

2026-06-10 & IIRu5hII & Aufgaben &

  Eintrag 18: Android-Projekt geöffnet, Compose-Grundstruktur geprüft,
  \texttt{MainActivity}/\texttt{SafetySpotApp} vorbereitet.
  Erste Orientierung an Package-Struktur \texttt{ui/auth}, \texttt{ui/home},
  \texttt{ui/scenarios}, \texttt{ui/ranking}, \texttt{ui/profile}.
  \newline\textit{Ergebnis:} Projekt kann als mobile Compose-App weiterentwickelt werden. \\

2026-06-10 & IIRu5hII & Ergebnis &

  Eintrag 19: Grundlegendes Design-System aufgebaut: Markenfarben,
  App-Hintergrund, Kartenfarben, Kategorie-Akzente und Schwierigkeitsfarben.
  \newline\textit{Ergebnis:} \texttt{Color.kt} enthält BrandGreen, BrandBlue, PointsYellow,
  ChemieBlueTint, WerkraumOrange, SportGreen, TechnikPurple, TrafficRed. \\

2026-06-11 & IIRu5hII & Aufgaben &

  Eintrag 20: Navigationskonzept umgesetzt: Auth als Startscreen,
  Hauptnavigation mit Home/Szenarien/Ranking/Profil sowie
  Detailrouten für Szenario-Detail, Gameplay und Profil-Unterseiten.
  \newline\textit{Ergebnis:} \texttt{SafetySpotApp} enthält NavHost und BottomBar-Logik. \\

2026-06-11 & IIRu5hII & Ergebnis &

  Eintrag 21: Bottom Navigation aufgebaut. Leiste wird nur auf Hauptscreens
  angezeigt, auf Detail-/Spielseiten ausgeblendet.
  \newline\textit{Ergebnis:} Navigation entspricht Mockups. \\

2026-06-11 & IIRu5hII & Ergebnis &

  Eintrag 22: AuthScreen erstellt. Enthält Branding, Login-/Registrieren-/
  Gastzugang-Logik als UI-Prototyp; Einstieg in die App.
  \newline\textit{Ergebnis:} Frontend ohne echte Zugangsdaten testbar. \\

11.06.2026 & 0qln & Ergebnis &

  Eintrag 23: Projekt-weite Struktur-Überarbeitung (PR \#28, \texttt{9ee221e}).
  Anpassungen an Build-Skripten und Projekt-Struktur nach erstem
  Feedback und paralleler Backend-Entwicklung.
  \newline\textit{Ergebnis:} Konsistente Projektstruktur über alle drei Module. \\

12.06.2026 & D0rag3on & Versuchsaufbau &

  Eintrag 24: Anlegen der Paket- und Klassenstruktur für globales Error-Handling.
  Neue Pakete: \texttt{dto}, \texttt{exception}.
  Erstellt als initiale Stubs:
  \texttt{ErrorResponse} (einheitliches DTO),
  \texttt{AccessDeniedException},
  \texttt{EntityNotFoundException},
  \texttt{GlobalExceptionHandler}.
  Alle Klassen bewusst leer gehalten – Struktur vor Implementierung. \\

2026-06-12 & IIRu5hII & Ergebnis &

  Eintrag 25: HomeScreen aufgebaut: Begrüßung, Level, Punkte, Streak,
  Weiterlernen-Bereich und Kategorien.
  \newline\textit{Ergebnis:} Zentraler Einstieg für Schüler nach Login. \\

2026-06-12 & IIRu5hII & Ergebnis &

  Eintrag 26: Wiederverwendbare Komponenten begonnen:
  \texttt{MetricCard}, \texttt{StreakRow}, \texttt{SafetyProgressBar},
  \texttt{SectionHeader}, \texttt{CategoryTile}, \texttt{ScenarioCard}.
  \newline\textit{Ergebnis:} Reduziert Dopplungen; spätere Anpassungen einfacher. \\

2026-06-12 & IIRu5hII & Idee &

  Eintrag 27: Für erste UI-Version bewusst Mockdaten genutzt,
  da Backend noch nicht vollständig angebunden.
  \newline\textit{Ergebnis:} Mockdaten werden schrittweise durch API-Daten ersetzt. \\

2026-06-13 & IIRu5hII & Ergebnis &

  Eintrag 28: ScenariosScreen: Titel, Suchfeld, Filtertabs
  (Alle/Neu/Beliebt/Abgeschlossen) und Szenario-Kartenliste.
  \newline\textit{Ergebnis:} Nutzer kann Szenarien visuell auswählen und filtern. \\

2026-06-13 & IIRu5hII & Ergebnis &

  Eintrag 29: \texttt{sampleScenarios} angelegt:
  Chemieraum, Werkraum, Sportunterricht, Straßenverkehr, Technikraum.
  \newline\textit{Ergebnis:} Entspricht Kategorien aus Projektvision und Mockups. \\

2026-06-13 & IIRu5hII & Problem &

  Eintrag 30: Umlaute und Sonderzeichen können in Code/Assets
  zu Darstellungsproblemen führen.
  \newline\textit{Ergebnis:} Teilweise ASCII-Schreibweisen in Mockdaten verwendet. \\

14.06.2026 & D0rag3on & Messergebnis &

  Eintrag 31: Vollständige Implementierung des \texttt{GlobalExceptionHandler}
  (45 Zeilen, \texttt{@RestControllerAdvice}).
  Alle definierten Ausnahmetypen werden auf HTTP-Statuscodes gemappt und als
  \texttt{ErrorResponse}-DTO zurückgegeben.
  Das globale Error-Handling ist damit abgeschlossen und einsatzbereit. \\

2026-06-14 & IIRu5hII & Ergebnis &

  Eintrag 32: \texttt{ScenarioDetailScreen} erstellt: Kategorie, Titel,
  Schwierigkeit, Aufgabenanzahl und Status vor dem Spielstart.
  \newline\textit{Ergebnis:} Zwischenschritt zwischen Szenarienliste und Gameplay. \\

2026-06-14 & IIRu5hII & Ergebnis &

  Eintrag 33: \texttt{ScenarioPlayScreen} aufgebaut: Fortschrittsanzeige,
  Punkteanzeige, Aufgabeninhalt, Kategorie-/Mascot-Bereich und
  große Antwortbuttons (gefährlich / nicht gefährlich).
  \newline\textit{Ergebnis:} Kernfunktion der App im Prototyp spielbar. \\

2026-06-14 & IIRu5hII & Annahmen/Hypothese &

  Eintrag 34: Hypothese: Zwei große Antwortbuttons sind für Schulkinder
  verständlicher als komplexes Multiple-Choice.
  \newline\textit{Ergebnis:} True/False-Modus für Version 1 priorisiert. \\

15.06.2026 & D0rag3on & Versuchsaufbau &

  Eintrag 35: Anlegen des vollständigen Projektskeletts für \texttt{ssbackend}:
  25 Java-Klassen über alle fachlichen Pakete verteilt.
  Pakete: \texttt{attempt} (8 Klassen), \texttt{auth} (2),
  \texttt{enums} (3), \texttt{leaderboard} (3),
  \texttt{progress} (2), \texttt{school} (3), \texttt{user} (4).
  Alle Klassen als Stubs; Ziel: Gesamtstruktur des Backends abbilden,
  bevor Implementierung beginnt. \\

2026-06-15 & IIRu5hII & Ergebnis &

  Eintrag 36: Feedbacklogik im Spielscreen ergänzt.
  Nach einer Antwort: korrekt/falsch hervorgehoben, Feedbacktext angezeigt.
  \newline\textit{Ergebnis:} Direktes Feedback unterstützt Lerneffekt. \\

2026-06-15 & IIRu5hII & Ergebnis &

  Eintrag 37: Lokaler Punktestand und Abschlusslogik eingeführt.
  Nach Abschluss eines Szenarios: Punkte addiert, Szenario als abgeschlossen markiert.
  \newline\textit{Ergebnis:} Gamification-Effekt ohne Backend-Persistenz. \\

2026-06-15 & IIRu5hII & Bewertung &

  Eintrag 38: Mock-first-Ansatz bewertet: Vorteil: schneller UI-Fortschritt;
  Nachteil: spätere Umstellung auf echte API-Modelle.
  \newline\textit{Ergebnis:} Entscheidung bleibt sinnvoll; Prototyp früh prüfbar. \\

2026-06-16 & IIRu5hII & Ergebnis &

  Eintrag 39: RankingScreen: Scope-Tabs (Klasse/Schule/Weltweit),
  Top-3-Darstellung, Liste weiterer Spieler mit Hervorhebung des aktuellen Nutzers.
  \newline\textit{Ergebnis:} Gamification-Ziel Punkte/Ranking sichtbar. \\

2026-06-16 & IIRu5hII & Ergebnis &

  Eintrag 40: \texttt{sampleLeaderboard} angelegt und dynamisch mit aktuellem
  Nutzer kombiniert. Lokaler Punktestand kann den Rang beeinflussen.
  \newline\textit{Ergebnis:} Ranking reagiert im Prototyp auf abgeschlossene Szenarien. \\

2026-06-16 & IIRu5hII & Risiko/Beobachtung &

  Eintrag 41: Leaderboard kann motivieren, aber auch Druck erzeugen.
  \newline\textit{Ergebnis:} Eigener Nutzer positiv hervorgehoben;
  spätere Option: Klassenranking abschaltbar. \\

2026-06-17 & IIRu5hII & Ergebnis &

  Eintrag 42: ProfileScreen erstellt: Avatar, Name, Level, XP-Balken,
  Punkte, Streak, Abzeichen und Menügruppen.
  \newline\textit{Ergebnis:} Profil macht Fortschritt und Gamification persönlich sichtbar. \\

2026-06-17 & IIRu5hII & Ergebnis &

  Eintrag 43: Profil-Unterseiten ergänzt:
  Mein Fortschritt, Abzeichen und Meine Szenarien.
  \newline\textit{Ergebnis:} Nutzer kann Szenarien und Auszeichnungen nachvollziehen. \\

2026-06-17 & IIRu5hII & Ergebnis &

  Eintrag 44: \texttt{SettingsScreen} und \texttt{HelpFeedbackScreen} ergänzt.
  \newline\textit{Ergebnis:} Erhöht Eindruck einer vollständigen App. \\

2026-06-18 & IIRu5hII & Aufgaben &

  Eintrag 45: Mockup-Screens mit realem Compose-Layout abgeglichen.
  Abstände, Karten, Rundungen und Farben angepasst.
  \newline\textit{Ergebnis:} UI wirkt wie Projekt-Mockups, nicht wie Standard-Demo. \\

2026-06-18 & IIRu5hII & Ergebnis &

  Eintrag 46: Szenario-/Kategoriedaten verfeinert.
  Chemie, Werkraum, Sport, Technik und Verkehr als Lernbereiche vorhanden.
  \newline\textit{Ergebnis:} Struktur lässt sich durch Backend-Kategorien erweitern. \\

2026-06-18 & IIRu5hII & Idee &

  Eintrag 47: SafetySpot-Drache/Maskottchen als unterstützendes Element vorgesehen.
  Für Feedback, Motivation und kindgerechte Ansprache.
  \newline\textit{Ergebnis:} Mascot-Bereich im Spielscreen vorgesehen; echte Assets später. \\

18.06.2026 & 0qln & Aufgaben &

  Eintrag 48: Mobile-Spezifikations-Issues ISSUE-012 bis ISSUE-022 erstellt
  (Commits \texttt{2161ee6}–\texttt{701aa4c}).
  Themen: Foundation/Build-Config, Design-System, Networking/Auth-Infrastruktur,
  Navigation-Shell, Auth-Screen, Home-Screen, Scenarios-Browse,
  Szenario-Gameplay, Ranking-Screen, Profil-Screen, Offline-Cache-Sync.
  \newline\textit{Ergebnis:} Vollständige Spezifikation der mobilen App-Screens
  als Grundlage für IIRu5hII's Frontend-Entwicklung und eigene Clickdummy-Arbeit. \\

2026-06-19 & IIRu5hII & Ergebnis &

  Eintrag 49: UI-Komponenten weiter standardisiert:
  einheitliche Cards, Chips, ProgressBars, Textfarben und Buttons.
  \newline\textit{Ergebnis:} Verbessert Wartbarkeit und konsistente Darstellung. \\

2026-06-19 & IIRu5hII & Problem &

  Eintrag 50: Durch lokale UI-Zustände drohten inkonsistente Punkte/
  abgeschlossene Szenarien beim Navigieren.
  \newline\textit{Ergebnis:} Zustand in \texttt{SafetySpotApp} zentralisiert:
  \texttt{sessionPoints} und \texttt{completedScenarioIds}. \\

2026-06-19 & IIRu5hII & Lösung &

  Eintrag 51: \texttt{completedScenarioIds} als \texttt{mutableStateListOf},
  \texttt{sessionPoints} als Compose-State.
  \newline\textit{Ergebnis:} Lokaler Prototyp ohne Backend trotzdem konsistent. \\

2026-06-20 & IIRu5hII & Ergebnis &

  Eintrag 52: Navigation zwischen Szenario-Liste, Detailansicht und Gameplay
  geprüft und verbessert.
  \newline\textit{Ergebnis:} Back-Navigation und Starten eines Szenarios funktionieren. \\

2026-06-20 & IIRu5hII & Test/Beobachtung &

  Eintrag 53: Manueller Test des Kernflows:
  Auth → Home → Szenarien → Detail → Spielen → Ergebnis → Ranking/Profil.
  \newline\textit{Ergebnis:} Kompletter Schüler-Flow im Frontend testbar. \\

2026-06-20 & IIRu5hII & Offene Punkte &

  Eintrag 54: Persistenz fehlt: Punktestand und Abschlussstatus gehen
  nach App-Neustart verloren.
  \newline\textit{Ergebnis:} Für nächste Stufe: Backend/Room/DataStore. \\

2026-06-21 & IIRu5hII & Ergebnis &

  Eintrag 55: Profil- und Rankingdaten mit lokalen Spielständen verknüpft.
  \newline\textit{Ergebnis:} Gamification fühlt sich zusammenhängender an. \\

2026-06-21 & IIRu5hII & Bewertung &

  Eintrag 56: Frontend-Prototyp erfüllt wichtigste UI-Anforderungen:
  simpel, mobil, visuell freundlich, große Bedienelemente, spielerischer Ablauf.
  \newline\textit{Ergebnis:} Basis für Demo und spätere Backend-Integration vorhanden. \\

2026-06-21 & IIRu5hII & Aufgaben &

  Eintrag 57: Code bereinigt und in UI-Pakete aufgeteilt:
  auth, home, scenarios, scenario, ranking, profile, navigation, theme, components.
  \newline\textit{Ergebnis:} Ordnung erleichtert Teamarbeit. \\

21.06.2026 & 0qln & Aufgaben &

  Eintrag 58: Frontend-Planung dokumentiert (PR \#43, \texttt{d01cb56}).
  Entscheidungen zu UI-Architektur, Screen-Übergängen und Datenfluss
  zwischen Frontend-Screens festgehalten.
  \newline\textit{Ergebnis:} Gemeinsames Verständnis für Frontend-Strukturierung
  mit IIRu5hII. \\

22.06.2026 & D0rag3on & Ergebnis &

  Eintrag 59: JWT-Authentifizierung vollständig implementiert (PR \#31):
  \texttt{AuthController} (Login, Logout, Register), \texttt{AuthService},
  \texttt{JwtUtil}, \texttt{JwtAuthFilter}, \texttt{AuthToken} (Persistenz),
  \texttt{AuthTokenRepository}, \texttt{SecurityConfig} (stateless JWT),
  \texttt{SecurityUserDetailsService}.
  Außerdem: Foreign-Key-Beziehungen zwischen \texttt{User}, \texttt{School}
  und \texttt{SchoolClass} korrekt verknüpft. \\

22.06.2026 & D0rag3on & Notiz &

  Eintrag 60: JWT-Secret aus \texttt{application.properties} als
  Umgebungsvariable externalisiert (PR \#45).
  Der hartcodierte Wert wird jetzt aus \texttt{\$\{JWT\_SECRET\}} gelesen;
  kein sensitives Secret verbleibt im Repository.
  Entspricht dem definierten Sicherheitskonzept.
  Commit erscheint zweimal in der Historie (Cherry-Pick nach Merge-Verlust). \\

2026-06-22 & IIRu5hII & Recherche &

  Eintrag 61: Abgleich mit Backend-/Informationsmodell.
  Benötigte API-Daten: User/Profile, Kategorien, Szenarien, Fortschritt,
  Leaderboard, Attempt/Answers.
  \newline\textit{Ergebnis:} Vorbereitung auf Entfernung harter Mockdaten. \\

2026-06-22 & IIRu5hII & Aufgaben &

  Eintrag 62: Mögliche DTO-/Repository-Grenzen identifiziert.
  UI soll langfristig über Repository/ViewModel-Daten arbeiten.
  \newline\textit{Ergebnis:} MVVM/Repository-Struktur als Zielarchitektur festgehalten. \\

2026-06-22 & IIRu5hII & Offene Punkte &

  Eintrag 63: Mobile ViewModels und echte Repository-Implementierungen
  noch nicht vollständig vorhanden.
  \newline\textit{Ergebnis:} Für späteres Refactoring: UI-State-Klassen und Tests. \\

22.06.2026 & 0qln & Versuchsaufbau &

  Eintrag 64: Kern-JPA-Entitäten implementiert (PR \#44, \texttt{f5af1c2}).
  Refactoring der Stub-Entitäten von D0rag3on's Basis:
  \texttt{User}, \texttt{School}, \texttt{SchoolClass}, \texttt{Scenario},
  \texttt{Category}, \texttt{Attempt} etc.\ mit Lombok-Annotations
  (\texttt{@Data}, \texttt{@Builder}, \texttt{@Entity}).
  Hibernate konfiguriert (\texttt{spring.jpa.hibernate.ddl-auto=update}).
  Erste Tests zum Laufen gebracht.
  \newline\textit{Notiz:} PR \#44 wartete explizit auf PR \#31 (D0rag3on's Auth):
  ,,Waiting for \#31 to complete''. \\

22.06.2026 & 0qln & Ergebnis &

  Eintrag 65: User-Management-API implementiert (PR \#46, \texttt{d8b7aa7}).
  Neue Klassen: \texttt{UserController} (CRUD-Endpoints \texttt{/api/v1/users}),
  \texttt{UserService}/\texttt{UserServiceImpl},
  \texttt{UserRepository} (Spring Data JPA),
  DTOs: \texttt{CreateUserRequest}, \texttt{UpdateUserRequest},
  \texttt{ResetPasswordRequest}, \texttt{UserResponse}.
  Passwörter werden stets gehasht (BCrypt).
  \newline\textit{Ergebnis:} Vollständige CRUD-API für Benutzer; Lehrer können
  Passwörter zurücksetzen. \\

22.06.2026 & 0qln & Ergebnis &

  Eintrag 66: Image- und ImageTag-Entitäten erstellt (PR \#47, \texttt{2c9afa2}).
  \texttt{Image}-Entität mit Metadaten (Beschreibung, Dateiname, Dateipfad),
  \texttt{ImageTag}-Entität (Verknüpfung Bild ↔ Nutzer ↔ Bewertung),
  \texttt{TagValue}-Enum (\texttt{DANGEROUS}/\texttt{SAFE}).
  \newline\textit{Ergebnis:} Datenmodell für kollaboratives Bild-Tagging im Backend. \\

22.06.2026 & 0qln & Ergebnis &

  Eintrag 67: Image-Management-API implementiert (PR \#48, \texttt{7a6cd2f}).
  \texttt{ImageController} (Upload/Download/CRUD \texttt{/api/v1/images}),
  \texttt{ImageService}/\texttt{ImageServiceImpl},
  \texttt{ImageDataService}/\texttt{ImageDataServiceImpl}
  (Datei-basierte Speicherung in \texttt{data/images/}),
  \texttt{ImageRepository}, \texttt{ImageTagRepository}.
  DTOs: \texttt{CreateImageRequest}, \texttt{UpdateImageRequest},
  \texttt{ImageResponse}, \texttt{ImageTagResultResponse}.
  \newline\textit{Ergebnis:} Bilder können hochgeladen, abgerufen und verwaltet werden. \\

22.06.2026 & 0qln & Ergebnis &

  Eintrag 68: Bild-Tagging-API implementiert (PR \#49, \texttt{92ab258}).
  \texttt{TagController} (\texttt{/api/v1/tags}), \texttt{TagService}/
  \texttt{TagServiceImpl}.
  Schüler können Bilder als \texttt{DANGEROUS} oder \texttt{SAFE} markieren;
  Duplikat-Tags werden mit \texttt{DuplicateTagException} abgelehnt.
  DTOs: \texttt{SubmitTagRequest}, \texttt{TagResponse},
  \texttt{ImageStatsResponse}.
  \newline\textit{Ergebnis:} Kern-Spielmechanik (Bild bewerten) im Backend vollständig. \\

22.06.2026 & 0qln & Versuchsaufbau &

  Eintrag 69: Nix Dev-Umgebung eingerichtet (PR \#50, \texttt{3d72b56}).
  \texttt{devenv.nix}: MariaDB, Java 25, Gradle deklarativ definiert;
  \texttt{devenv.yaml} und \texttt{devenv.lock} für Reproduzierbarkeit.
  \newline\textit{Ergebnis:} Einheitliche, reproduzierbare Entwicklungsumgebung
  für alle Teammitglieder ohne manuelle Tool-Installation. \\

22.06.2026 & 0qln & Ergebnis &

  Eintrag 70: Schülerfortschritt- und Lehrer-Stats-API implementiert
  (PR \#51, \texttt{5264857}).
  \texttt{Progress}-Entität (Schüler ↔ Szenario: Score, Attempts, Completed),
  \texttt{ProgressController}/\texttt{ProgressService}/\texttt{ProgressServiceImpl}
  (\texttt{/api/v1/progress}).
  \texttt{StatsController}/\texttt{StatsService}/\texttt{StatsServiceImpl}
  (\texttt{/api/v1/stats}): Klassen-Statistiken pro Lehrer.
  DTOs: \texttt{ProgressEntryResponse}, \texttt{ProgressSummaryResponse},
  \texttt{ClassStatsResponse}, \texttt{StudentStatEntry}.
  \newline\textit{Ergebnis:} Lehrer können Klassenergebnisse einsehen;
  Schüler-Fortschritt wird serverseitig persistiert. \\

22.06.2026 & 0qln & Ergebnis &

  Eintrag 71: Build- \& Release-Pipeline aufgebaut (PR \#52, \texttt{95f8be5}).
  \texttt{Dockerfile} für ssbackend,
  GitHub-Actions-Workflow \texttt{azure-container-webapp.yml}:
  Docker-Build → Push zu Azure Container Registry → Deploy zu Azure App Service.
  \texttt{application-prod.properties} mit Azure-SQL-Datasource-Konfiguration.
  \newline\textit{Ergebnis:} Automatisierter Deployment-Prozess vom Commit
  bis zur Produktionsumgebung auf Azure. \\

22.06.2026 & 0qln & Problem &

  Eintrag 72: Azure-Container-Webapp-Workflow musste direkt nach Anlegen
  zweimal korrigiert werden (Commits \texttt{3ae7804}, \texttt{a9b36ee}).
  Problem: Fehlende Umgebungsvariablen und inkorrekte Workflow-Trigger.
  \newline\textit{Ergebnis:} Workflow läuft stabil nach zwei Korrekturen. \\

2026-06-23 & IIRu5hII & Ergebnis &

  Eintrag 73: Frontend-Prototyp vor Git-Historie konsolidiert.
  Wichtigste Screens mit Mockdaten verbunden.
  \newline\textit{Ergebnis:} Grundlage für anschließende Integration-Commits. \\

2026-06-23 & IIRu5hII & Problem &

  Eintrag 74: Frühe Arbeiten ohne Commits — technische Historie nicht
  vollständig aus Git rekonstruierbar.
  \newline\textit{Ergebnis:} Laborbuch rekonstruiert frühe Schritte aus Quellcode und Mockups. \\

2026-06-23 & IIRu5hII & Offene Punkte &

  Eintrag 75: Backend-Anbindung, echte Auth, Bilddaten und persistenter
  Fortschritt müssen angebunden werden.
  \newline\textit{Ergebnis:} Nächster Fokus: API/Retrofit und Integration. \\

23.06.2026 & 0qln & Problem &

  Eintrag 76: Spring-Boot-Tests liefen nicht (Commit-Kontext PR \#53,
  \texttt{3cc8d2a}).
  Ursachen: Fehlende H2-Testdatenbank-Konfiguration, Spring-Security-
  Testkonfiguration nicht gesetzt, Testfixtures unvollständig.
  \newline\textit{Lösung:} H2 In-Memory für Testprofil konfiguriert,
  \texttt{@SpringBootTest}-Annotations ergänzt, 17 Testklassen lauffähig. \\

23.06.2026 & 0qln & Ergebnis &

  Eintrag 77: CI-Test-Workflow erstellt (PR \#56, \texttt{63756d9}).
  \texttt{.github/workflows/test.yml}: Gradle-Tests auf jedem
  Push und Pull-Request automatisch ausgeführt.
  \newline\textit{Ergebnis:} Kontinuierliche Qualitätssicherung;
  Regressions-Erkennung bei jedem Commit. \\

23.06.2026 & 0qln & Ergebnis &

  Eintrag 78: Mobile-UI-Clickdummy vollständig implementiert (PR \#54,
  \texttt{f9ab2db}). Ca.\ 4\,072 Zeilen Kotlin/Compose neu hinzugefügt,
  27 Dateien geändert.
  Screens: \texttt{AuthScreen} (Login/Register), \texttt{HomeScreen}
  (Level, Punkte, Streak, Kategorien), \texttt{ScenariosScreen} (Filter, Suche),
  \texttt{ScenarioDetailScreen}, \texttt{ScenarioPlayScreen} (Gameplay,
  Antwort-Buttons, Feedback), \texttt{RankingScreen} (Scope-Tabs, Top-3),
  \texttt{ProfileScreen} mit Unterseiten.
  Navigation via NavHost/BottomBar, Compose-Theme mit SafetySpot-Farben,
  Mock-Daten für vollständige Offline-Durchläufe.
  \newline\textit{Ergebnis:} Vollständig klickbarer Prototyp zum Demo-Termin bereit. \\

2026-06-24 & IIRu5hII & Ergebnis &

  Eintrag 79: Commit \texttt{7d86b5d}: Retrofit kann nun auch Plain Text lesen.
  Retrofit Scalars, DTOs, Auth/Login, Backend-Bilder sowie Home/Szenarien
  aus Backend vorbereitet und angebunden.
  \newline\textit{Ergebnis:} Beginn der konkreten Mobile-Backend-Integration. \\

2026-06-24 & IIRu5hII & Aufgaben &

  Eintrag 80: DTOs und API-Kommunikation für Login/Auth und Datenabruf geprüft.
  \newline\textit{Ergebnis:} Vorbereitung auf Zusammenspiel ssmobile ↔ ssbackend. \\

2026-06-24 & IIRu5hII & Ergebnis &

  Eintrag 81: Backend-Bilder als Quelle für Szenario-/Aufgabenbilder berücksichtigt.
  \newline\textit{Ergebnis:} Wichtig für spätere Hotspot-Gefahrensituationen. \\

2026-06-24 & IIRu5hII & Ergebnis &

  Eintrag 82: Home- und Szenarienbereiche auf spätere Backend-Daten vorbereitet.
  \newline\textit{Ergebnis:} Mockdaten bleiben als Fallback/Demo nützlich. \\

2026-06-24 & IIRu5hII & Ergebnis &

  Eintrag 83: Commit \texttt{09beaca}: Merge von origin/main in
  frontend-functionality. Übernommen: Test-Workflow, README,
  DemoDataService, AdminShellCommands.
  \newline\textit{Ergebnis:} Frontend-Branch mit aktuellem main synchronisiert. \\

2026-06-24 & IIRu5hII & Problem &

  Eintrag 84: Merges können Konflikte zwischen Frontend-Prototyp und
  Backend-/DemoDataService-Änderungen verursachen.
  \newline\textit{Ergebnis:} Nach Merge geprüft, dass relevante Änderungen erhalten bleiben. \\

2026-06-24 & IIRu5hII & Bewertung &

  Eintrag 85: CI/Test-Workflow und AdminShellCommands wichtig für Teamarbeit.
  \newline\textit{Ergebnis:} Projekt reproduzierbarer und wartbarer. \\

24.06.2026 & 0qln & Ergebnis &

  Eintrag 86: Admin-Shell und Demo-Daten implementiert (PR \#58,
  \texttt{8dc1fe1}).
  \texttt{AdminShellCommands} (Spring Shell): interaktive CLI-Kommandos
  (\texttt{create-school}, \texttt{create-class}, \texttt{create-user},
  \texttt{seed}, \texttt{list-users} etc.).
  \texttt{DemoDataService}: automatischer Seed beim App-Start
  (Schule, Klassen, Schüler, Lehrer, Bilder, Tags).
  \newline\textit{Ergebnis:} Backend kann ohne manuelle SQL-Statements
  mit realistischen Testdaten befüllt werden. \\

24.06.2026 & 0qln & Notiz &

  Eintrag 87: Temporärer Workaround in ssmobile eingeführt (Commit
  \texttt{221ef31}) und am 27.06.2026 wieder rückgängig gemacht
  (Commit \texttt{7a80d24}).
  Hintergrund: Schnelle Umgehung eines Kompatibilitätsproblems zwischen
  Frontend-Functionality-Branch und Backend-Änderungen während
  aktiver Merge-Phase.
  \newline\textit{Ergebnis:} Nur temporär; saubere Lösung nach Merge eingespielt. \\

2026-06-25 & IIRu5hII & Ergebnis &

  Eintrag 88: Commit \texttt{5a24e56}: Merge branch main into frontend-functionality.
  DemoDataService-Änderungen übernommen.
  \newline\textit{Ergebnis:} Datenbasis/Seed-Daten mit Frontend-Branch kompatibel. \\

2026-06-25 & IIRu5hII & Ergebnis &

  Eintrag 89: Commit \texttt{d804323}: weiterer Merge von main in
  frontend-functionality. Backend-/Security-Konfiguration übernommen.
  \newline\textit{Ergebnis:} Auth-relevante Konfiguration im Branch vorhanden. \\

2026-06-25 & IIRu5hII & Aufgaben &

  Eintrag 90: Nach Merges geprüft, ob Frontend-Navigation, Demo-Daten und
  Backend-Konfiguration zusammenpassen.
  \newline\textit{Ergebnis:} Keine Regression im klickbaren Prototyp. \\

2026-06-25 & IIRu5hII & Ergebnis &

  Eintrag 91: Commit \texttt{41832d1}: Kategorie-Icons, Punkte-/Suchicon,
  Weitermachen-Logik, lokaler Fortschritt sowie Profilname/Avatar-Fallback ergänzt.
  \newline\textit{Ergebnis:} UI wirkt vollständiger und robuster. \\

2026-06-25 & IIRu5hII & Lösung &

  Eintrag 92: Fallbacks für Profilname und Avatar eingeführt.
  \newline\textit{Ergebnis:} App bleibt bei fehlenden Backend-Daten nutzbar. \\

2026-06-25 & IIRu5hII & Ergebnis &

  Eintrag 93: Weitermachen-Logik verbessert: Home kann begonnenes Szenario
  anzeigen und Nutzer zurück in den Lernfluss führen.
  \newline\textit{Ergebnis:} Erhöht Nutzerbindung und passt zur Gamification-Idee. \\

2026-06-25 & IIRu5hII & Test/Beobachtung &

  Eintrag 94: Nach Icon- und Fortschrittsänderungen wurden Home, Szenarien,
  Profil und Ranking manuell geprüft.
  \newline\textit{Ergebnis:} Keine offensichtlichen Brüche im Hauptfluss. \\

25.06.2026 & 0qln & Ergebnis &

  Eintrag 95: Demo-Daten korrigiert (PR \#59, \texttt{e766cb6}).
  Fehler im \texttt{DemoDataService}: falsche Beziehungs-IDs und
  inkonsistente Enum-Werte behoben.
  \newline\textit{Ergebnis:} Backend startet mit korrekten und konsistenten
  Testdaten. \\

25.06.2026 & 0qln & Ergebnis &

  Eintrag 96: H2-Konsole für Dev-Profil aktiviert (PR \#60, \texttt{c24bcdf}).
  \texttt{spring.h2.console.enabled=true} in
  \texttt{application-dev.properties}.
  \newline\textit{Ergebnis:} Direkter Datenbankzugriff über Browser während
  Entwicklung erleichtert Debugging. \\

25.06.2026 & 0qln & Ergebnis &

  Eintrag 97: Dateibasierte Bildspeicherung implementiert (PR \#61,
  \texttt{659fa1a}).
  \texttt{ImageDataServiceImpl} schreibt und liest Binärdaten unter
  \texttt{ssbackend/data/images/}.
  Verzeichnis mit gitignore-Ausnahme für Demo-Daten.
  \newline\textit{Ergebnis:} Bilder werden persistent auf dem Dateisystem
  gespeichert und per REST-Endpunkt ausgeliefert. \\

25.06.2026 & 0qln & Aufgaben &

  Eintrag 98: Debugging-Session Backend (Commit \texttt{1d41833}).
  Untersuchung von Fehlern bei der Bild-Upload/Download-Kette zwischen
  ssmobile und ssbackend.
  \newline\textit{Ergebnis:} Probleme identifiziert und in Folgecommits behoben. \\

25.06.2026 & 0qln & Ergebnis &

  Eintrag 99: Szenario-Demo-Bilder hinzugefügt (Commit \texttt{60f19a1}).
  10 PNG-Bilder (\texttt{scenario-00.png} bis \texttt{scenario-09.png})
  unter \texttt{ssbackend/src/main/resources/demo/images/}.
  Diese werden vom \texttt{DemoDataService} beim Seed in die Datenbank geladen.
  \newline\textit{Ergebnis:} Vollständiger Demo-Betrieb mit realistischen
  Bild-Szenarien möglich; IIRu5hII kann Backend-Bilder direkt testen. \\

26.06.2026 & D0rag3on & Ergebnis &

  Eintrag 100: School-/Klassen-Logik vollständig implementiert (PR \#55):
  \texttt{SchoolClass}-Entität mit JPA-Beziehung zu \texttt{School};
  \texttt{SchoolClassController}, \texttt{SchoolClassService},
  \texttt{SchoolClassRepository}.
  Ergänzend: \texttt{SchoolController} und \texttt{SchoolService} fertiggestellt.
  DTOs: \texttt{CreateClass}, \texttt{CreateSchool}, \texttt{UpdateSchoolClass},
  \texttt{UpdateSchoolRequest}. \\

2026-06-26 & IIRu5hII & Ergebnis &

  Eintrag 101: Commit \texttt{edab15a}: restliche Icons ergänzt.
  Profil-Icons, Filter-Icon und Avatar-Icon hinzugefügt;
  Profil- und Szenario-UI angepasst.
  \newline\textit{Ergebnis:} Letzter Feinschliff am visuellen Eindruck. \\

2026-06-26 & IIRu5hII & Ergebnis &

  Eintrag 102: Profil-UI weiter an Mockups angepasst.
  Menüpunkte und Statuswerte visuell klarer erkennbar.
  \newline\textit{Ergebnis:} Profilseite für Schüler verständlicher und vollständiger. \\

2026-06-26 & IIRu5hII & Ergebnis &

  Eintrag 103: Szenario-UI angepasst. Filter- und Szenariodarstellung optisch verbessert.
  \newline\textit{Ergebnis:} Szenarienliste leichter zu scannen; entspricht App-Stil. \\

2026-06-26 & IIRu5hII & Bewertung &

  Eintrag 104: Die UI ist als funktionsfähiger Frontend-Prototyp vorzeigbar:
  Login/Gastzugang, Home, Szenarien, Spielscreen, Ranking und Profil.
  \newline\textit{Ergebnis:} Erfüllt MVP-Schwerpunkt aus Projektvision. \\

2026-06-26 & IIRu5hII & Offene Punkte &

  Eintrag 105: Echte Persistenz und vollständige Backend-Anbindung müssen
  weiter stabilisiert werden.
  \newline\textit{Ergebnis:} Nächste Arbeit: Repository/ViewModel-Schicht,
  Auth-State, Tokenhandling, Cache. \\

2026-06-26 & IIRu5hII & Offene Punkte &

  Eintrag 106: Lehrer-/Adminbereich im Mobile-Frontend noch nicht umgesetzt.
  \newline\textit{Ergebnis:} Für später: Klassen, Schülerverwaltung,
  Szenario-Editor, Statistiken. \\

2026-06-26 & IIRu5hII & Offene Punkte &

  Eintrag 107: Offline-Modus konzeptionell vorgesehen, aber nicht umgesetzt.
  \newline\textit{Ergebnis:} Mit Room/DataStore/Pending Attempts planen. \\

2026-06-26 & IIRu5hII & Schlussfolgerung &

  Eintrag 108: Mock-first-Ansatz hat sich bewährt: trotz fehlender früher Commits
  schnell nutzbarer Prototyp mit Design, Navigation und Lernflow.
  \newline\textit{Ergebnis:} Für finale Abgabe: Git-Commits künftig kleinschrittiger. \\

26.06.2026 & 0qln & Ergebnis &

  Eintrag 109: Projekt-weite Code-Formatierung angewendet (PR \#62,
  \texttt{3bda336}).
  Gradle-Format-Task über alle Java-Dateien in \texttt{ssbackend/}
  ausgeführt; einheitlicher Code-Stil durchgesetzt.
  \newline\textit{Ergebnis:} Diffs in Folge-PRs werden sauberer;
  Review-Aufwand sinkt. \\

27.06.2026 & 0qln & Ergebnis &

  Eintrag 110: Szenario-Bilder aktualisiert (Commit \texttt{4736baf}).
  Bilder-Set überarbeitet und mit dem aktuellen Stand des Backends
  abgeglichen.
  \newline\textit{Ergebnis:} Demo-Bilder im Frontend und Backend konsistent. \\

27.06.2026 & 0qln & Ergebnis &

  Eintrag 111: Android-APK-Release-Workflow erstellt (PR \#64/\#65,
  \texttt{fbaf527}, \texttt{a54247d}, \texttt{485ea21}).
  \texttt{android-apk-release.yml}: automatischer APK-Build bei Git-Tag;
  APK als GitHub-Release-Asset veröffentlicht.
  Zwei Korrekturen nötig (PR \#64 initiale Erstellung, PR \#65 Fixup).
  \newline\textit{Ergebnis:} Vollständige Release-Pipeline für ssmobile;
  APK bei jedem Tag-Push automatisch verfügbar. \\

27.06.2026 & 0qln & Notiz &

  Eintrag 112: Branch \texttt{frontend-functionality} nach umfangreichen
  Merges und Koordination mit IIRu5hII in main zusammengeführt (PR \#63,
  \texttt{0027401}).
  Der PR enthielt Commits beider Autoren (IIRu5hII's Icons/Weitermachen-Logik
  und 0qln's Debugging/Bilder/Workaround-Rücknahme).
  \newline\textit{Ergebnis:} Finaler Stand der mobilen Frontend-Funktionalität
  in main integriert; GitHub-Release \texttt{v0.0.0} erstellt. \\
\end{sslabortabelle}

\section{Zusammenfassung}

\textbf{D0rag3on:} Backend-Grundgerüst, Maven-zu-Gradle-Migration, JWT-Authentifizierung,
globales Error-Handling und School-/Klassen-Verwaltung.

\textbf{IIRu5hII:} Android-Frontend (Kotlin/Compose) – Design-System, Navigation,
alle Screens (Auth, Home, Szenarien, Gameplay, Ranking, Profil) sowie Backend-Integration
und UI-Feinschliff.

\textbf{0qln:} Projektinfrastruktur, vollständige Backend-API (User, Image, Tag,
Progress, Stats), CI/CD-Pipelines (Azure, Android), Mobile-Clickdummy und Demo-Daten.
